\section{Glosario}

\subsection{Descripcion}

Este glosario define los terminos base del proyecto para evitar ambiguedades entre documentos.

\subsection{Alcance}

Aplica a todo el repositorio.

\subsection{Definiciones}

\begin{longtable}{p{0.28\textwidth}p{0.64\textwidth}}
\toprule
\textbf{Termino} & \textbf{Definicion} \\
\midrule
Usuario pendiente & Persona registrada que aun no puede participar. \\
Usuario activo & Persona habilitada para ingresar picks y competir. \\
Usuario rechazado & Persona cuya activacion fue denegada durante el flujo de alta. \\
Usuario bloqueado & Persona cuya participacion fue revocada por accion administrativa. \\
Kickoff & Fecha y hora oficial de inicio del partido en el sistema. \\
Pick & Prediccion del usuario para un partido. \\
Marcador exacto & Prediccion de goles local y goles visita. \\
Outcome & Resultado del partido: gana local, empate o gana visita. \\
Lock & Momento a partir del cual un pick deja de ser editable. \\
Resultado oficial & Resultado utilizado para calcular puntajes. \\
Override manual & Correccion administrativa explicitamente auditada. \\
Leaderboard & Tabla de posiciones general o filtrada. \\
Fase KO & Cualquier fase eliminatoria posterior a grupos. \\
Archivado & Estado de cierre del torneo en modo solo lectura. \\
\bottomrule
\end{longtable}

\subsection{Reglas}

\begin{itemize}
  \item Un termino definido aqui debe usarse de forma consistente en el resto del repositorio.
  \item Si un termino cambia de significado, el glosario se actualiza primero.
\end{itemize}

\subsection{Edge Cases}

\begin{itemize}
  \item Si dos areas usan palabras distintas para el mismo concepto, prevalece el termino normalizado en este glosario.
\end{itemize}

\subsection{Invariantes}

\begin{itemize}
  \item No se usan sinonimos funcionales cuando puedan introducir ambiguedad.
\end{itemize}

\subsection{Preguntas abiertas}

\begin{itemize}
  \item Ninguna por ahora.
\end{itemize}
